\documentclass[11pt]{article}
\usepackage{fullpage,graphicx,hyperref,amsmath,amssymb}

\newcommand{\fig}[2]{\begin{figure}[h]\begin{center}%
  \includegraphics[scale=#2]{#1.pdf}\end{center}%
  \end{figure}}
\newcommand{\ceil}[1]{\left\lceil #1\right\rceil}

\begin{document}

\begin{center}\Large\bf 
CS 473 (Spring 2025)\\
Homework 8 (due Apr 17 Thu 10am)
\end{center}

\medskip\noindent
{\bf Instructions:} As in previous homeworks.

\begin{description}

\bigskip
\item[Problem 8.1:]\

\begin{enumerate}
\item[(a)] (80 pts)
Run the simplex method 
on the following linear program:
\[ \begin{array}{lrcl}
\mbox{maximize} & x_1 + x_2 &&\\
\mbox{s.t.} & 3x_1 + 2x_2 &\le& 27\\
&x_2&\le&8\\
& x_1+5x_2&\le&35\\
& 2x_1 + x_2 &\le& 17\\
&x_1,x_2&\ge& 0.
\end{array}\]
Start with the initial basic solution
$(\overline{x}_1,\overline{x}_2)=(0,0)$, and
{\em choose $x_1$ as the entering variable in the first iteration}.
Show the new slack form after every iteration.
\item[(b)] (20 pts)
Write down the dual of the linear program from (a).
What is the optimal dual solution?
\end{enumerate}


\bigskip
\item[Problem 8.2:]
\newcommand{\VEC}[1]{#1}%{\mathbf{#1}}
\newcommand{\R}{\mathbb{R}}
For two points $\VEC{p}=(p_1,\ldots,p_d)\in \R^d$ and 
$\VEC{q}=(q_1,\ldots,q_d)\in \R^d$,
their \emph{rectilinear distance} is defined as 
$$D(\VEC{p},\VEC{q})=|p_1-q_1|+\cdots+|p_d-q_d|.$$
In the {\em rectilinear 1-center} problem in $d$ dimensions,
we are given a set $P$ of $n$ points in $\R^d$,
we want to find a point $\VEC{q}\in \R^d$ (not necessarily in $P$)
that minimizes $\max_{\VEC{p}\in P} D(\VEC{p},\VEC{q})$.

\begin{enumerate}
\item[(a)] (30 pts)
First show that the problem can be solved in $O(n)$ time for
any constant dimension $d$.  How does the running time of your
algorithm grow as a function of $d$?

(Hint: find a small number of candidate points in $P$ that could be the 
farthest point from any $\VEC{q}$\ldots)
\item[(b)] (70 pts)
Show how to solve this problem for large (nonconstant) dimensions $d$ by
using linear programming.
The number of variables and constraints should be polynomial in $n$
and~$d$.  (Remember to justify correctness of your reduction.)
\end{enumerate}


\bigskip
\item[Problem 8.3:]
We are given $n$ tasks to perform.  Task $i$ requires
$p_i$ units of power consumption for a duration of $h_i$ hours.
At any moment in time, we can perform at most 3 different tasks, 
and at any moment in time, the total power consumption must be at most
$P$.  A task may be preempted (possibly multiple times) at
no extra cost.  The problem is to devise a schedule to
perform all $n$ tasks with the minimum total number of hours.

For example: for $n=5$ with $p_1=10$, $h_1=8.5$, $p_2=20$, $h_2=9$,
$p_3=60$, $h_3=4$, $p_4=80$, $h_4=3.5$, $p_5=90$, $h_5=2$, and $P=100$,
one feasible solution is to do tasks 1 and 5 for 2 hours,
then tasks 2 and 4 for $3.5$ hours,
then tasks 1, 2, and 3 for 4 hours,
then tasks 1 and 2 for $1.5$ hours,
and finally task 1 for 1 hour; the total number of hours is $12$.
(I did not check if this is optimal.  Also, for this small
example, the constraint that we can do at most 3 tasks at any time
is not important; but it could make a difference on larger instances.)

Describe how to solve this problem
using linear programming.
The number of variables and constraints should be polynomial in $n$.

\end{description}
\end{document}
